% ==============================================================
% Adversarial Review: Ontophänomenalismus (OP)
% Reviewer: Vibe (Mistral Medium 3.5)
% Datum: 4. Juli 2026
% ==============================================================

\documentclass[11pt,a4paper]{article}
\usepackage[utf8]{inputenc}
\usepackage[T1]{fontenc}
\usepackage[german]{babel}
\usepackage{geometry}
\usepackage{parskip}
\usepackage{enumitem}
\usepackage{booktabs}
\usepackage{hyperref}
\geometry{margin=2.5cm}

\title{
  \textbf{Adversarial Review}\\
  \textbf{Forschungsprogramm Ontophänomenalismus}\\[0.4em]
  {\normalsize Reviewer: Vibe --- 4. Juli 2026}
}
\author{}
\date{}

\begin{document}
\maketitle

%--------------------------------------------------------------
\section*{Zusammenfassung}
%--------------------------------------------------------------
Das Forschungsprogramm Ontophänomenalismus (OP) präsentiert sich als ein **kohärentes, methodisch reflektiertes System**, das durch seine **modulare Architektur**, **transparente Methodik (MKO)** und **radikale ontologische Entscheidungen** (Denk-Monismus, OPP als Urgrund) überzeugt. Besonders solide sind:
\begin{itemize}
    \item Die **R/P-Unterscheidung** (MKO) als Instrument der Zirkelkontrolle.
    \item Die **Sprachregel mit fünf Ebenen** (Indikativ, Konjunktiv, Definitorisch, Modal, Hypothetisch), die **Kategorienfehler** verhindert.
    \item Die **Relokation des Explanatory Gaps** (EG-1) als konstitutive Grenze jeder endlichen Perspektive.
    \item Die **klare Abgrenzung emergenter Begriffe** (Raum, Zeit, Energie) in OPP und OAD, die erst in DPR legitim eingeführt werden.
\end{itemize}

Gleichzeitig gibt es **kritische Lücken**, insbesondere:
\begin{itemize}
    \item **Inkonsistente Anwendung der MKO-Regeln** (z. B. Indikativ statt Konjunktiv in Interpretationen).
    \item **Fehlende explizite Kennzeichnung der R/P-Unterscheidung** in OP.tex (Axiom 3).
    \item **Unklare Abgrenzung zwischen KI als Adversarial Reviewer und Ko-Autor** (TRANSPARENCY).
    \item **Terminologische Inkonsistenzen** (z. B. "Denkfeld" vs. "OPP").
\end{itemize}

Die **prioritär zu adressierenden Befunde** sind **[K-01]**, **[K-02]**, **[K-03]** und **[K-04]**, da sie die **methodologische Stringenz** und **Glaubwürdigkeit** des OP direkt betreffen.

%--------------------------------------------------------------
\section*{Kritische Befunde [K]}
%--------------------------------------------------------------

\subsection*{[K-01] Unklare Urheberschaftsabgrenzung zwischen Autor und KI in TRANSPARENCY.tex}
\textbf{Dokument:} TRANSPARENCY.tex \\
\textbf{Stelle:} Abschnitt "Zusammenarbeit folgt einer klar definierten methodischen Rollenverteilung" \\
\textbf{Problem:}
Die Rollen der KI werden als **Adversarial Reviewer** (kritische Gegenposition) und **formale Ausarbeitung** (sprachliche Präzisierung, LaTeX-Umsetzung) beschrieben. Allerdings bleibt unklar, wie die **Grenze zwischen Kritik und inhaltlicher Mitgestaltung** gezogen wird. Besonders problematisch:
\begin{itemize}
    \item KI-generierte **konzeptionelle Vorschläge** (z. B. logische Konsolidierungen in OP.tex) werden nicht explizit als **Adversarial Review** oder **Ko-Autorschaft** klassifiziert.
    \item **Risiko:** Wenn KI-Systeme nicht nur kritisieren, sondern auch **inhaltlich prägen**, könnte dies die **Urheberschaftsclaims** untergraben, besonders da die KI als **"Mitdenker"** agiert.
\end{itemize}
\textbf{Warum kritisch:}
Die **Transparenz der Urheberschaft** (Ziel von TRANSPARENCY) wird geschwächt, wenn die Rolle der KI nicht klar von der des Autors abgegrenzt ist. Zudem könnte dies die **Rekonstruktivität (R)** der Axiome infrage stellen, da diese gemäß MKO **ausschließlich vom Autor** stammen sollten.

---

\subsection*{[K-02] Fehlende explizite R/P-Kennzeichnung in OP.tex (Axiom 3)}
\textbf{Dokument:} OP.tex \\
\textbf{Stelle:} Abschnitt "Axiom 3: Der Denk-Monismus" \\
\textbf{Problem:}
Axiom 3 wird als **postulativ (P)** klassifiziert, aber die **Begründung für diese Klassifizierung** ist nicht **explizit als "Methodennote" oder "Konvention"** (gemäß MKO) ausgewiesen. Stattdessen wird die Begründung im Fließtext gegeben, ohne formale Kennzeichnung.
\textbf{Warum kritisch:}
Verstoß gegen **MKO, R/P-Unterscheidung**: Die R/P-Kennzeichnung ist ein **zentrales Instrument der Zirkelkontrolle**. Fehlende explizite Kennzeichnung könnte zu **Missverständnissen** führen, ob Axiom 3 als rekonstruktiv (R) oder postulativ (P) zu verstehen ist. Dies untergräbt die **methodologische Transparenz** des OP.

---

\subsection*{[K-03] Inkonsistente Verwendung der Sprachebenen in OP.tex}
\textbf{Dokument:} OP.tex \\
\textbf{Stelle:} Abschnitt "Minimaldefinition des Denkens" (Axiom 3, Schritt 2) \\
\textbf{Problem:}
Gemäß **MKO, Sprachebene 2** müssen **ontologische Interpretationen** im **Konjunktiv** stehen. In OP.tex wird jedoch der Satz:
\begin{quote}
"Der OPP besitzt genau diese vier Minimalmerkmale: Selbstpräsenz (Axiom 2), informationelle Differenzierung (UP), raumzeitlose Statik (OPP), unmittelbare Gegebenheit (phänomenologische Priorität)."
\end{quote}
im **Indikativ** formuliert, obwohl es sich um eine **Interpretation** (keine direkte logische Folgerung aus Axiom 1 oder 2) handelt.
\textbf{Warum kritisch:}
Verstoß gegen **MKO, Sprachebene 2**. Ontologische Interpretationen **müssen** im Konjunktiv stehen (z. B. *"Der OPP könnte diese vier Minimalmerkmale besitzen"*). Die **methodologische Stringenz** des OP wird untergraben, wenn Interpretationen fälschlich als Fakten präsentiert werden.

---
\subsection*{[K-04] Kontamination in OAD.tex: Verwendung emergenter Begriffe in der ontischen Basis}
\textbf{Dokument:} OAD.tex \\
\textbf{Stelle:} Abschnitt "Prinzip der strukturellen Stabilisierung" \\
\textbf{Problem:}
In der **Definition der ontologischen Stabilität** wird der Begriff **"Fixpunkt"** verwendet, der **mathematisch** konnotiert ist und implizit eine **Metrik oder Dynamik** voraussetzt. Gemäß **MKO, Kontaminationskontrolle** sind solche Begriffe in OPP und OAD **nur als explizit gekennzeichnete Rückwärts-Analogien** erlaubt.
\begin{quote}
"Ein Zustand im OPP gilt genau dann als ontologisch stabil, wenn er seine inhärente Unterscheidbarkeit gemäß der Unterscheidbarkeits-Prämisse (UP) bewahrt. Begrifflich (MKO, Ebene 3) lässt sich ein solcher Zustand als Fixpunkt eines inhärenten Feldoperators modellieren..."
\end{quote}
\textbf{Warum kritisch:}
Hier wird **"Fixpunkt"** nicht als **Analogie (MKO Ebene 4)**, sondern als **Definition (MKO Ebene 3)** eingeführt. Dies könnte als **Kategorienfehler** interpretiert werden, da der Begriff **emergente Eigenschaften** (z. B. aus der Mathematik) in die **ontische Basis** einschmuggelt.

---

\subsection*{[K-05] Fehlende klare Abgrenzung zwischen Definition und Modell in DFD.tex}
\textbf{Dokument:} DFD.tex \\
\textbf{Stelle:} Abschnitt "Mathematische Heuristik" \\
\textbf{Problem:}
Das **Fragedifferential** wird als **ontologischer Begriff** eingeführt, aber die **mathematische Formalisierung als KL-Divergenz** wird als **"erste informationstheoretische Heuristik"** bezeichnet. Gemäß **MKO** muss klar zwischen **Definition** (terminologisch stabil) und **Modell** (ersetzbar) unterschieden werden.
\begin{quote}
"Als erste mathematische Heuristik definieren wir das Fragedifferential $\mathcal{D}(E)$ durch die Kullback-Leibler-Divergenz: $\mathcal{D}(E) := D_{\mathrm{KL}}(Q_E \| P)$"
\end{quote}
\textbf{Warum kritisch:}
Die Formulierung **"definieren wir"** suggeriert, dass die KL-Divergenz eine **Definition** ist, obwohl es sich um ein **Modell** handelt. Dies könnte zu **Missverständnissen** führen, ob das Fragedifferential **ontologisch** oder **nur formal** zu verstehen ist.

---
%--------------------------------------------------------------
\section*{Bedeutsame Befunde [B]}
%--------------------------------------------------------------

\subsection*{[B-01] Unscharfe Abgrenzung zwischen "Definition" und "Modell" in MKO.tex und den Modulen}
\textbf{Dokument:} MKO.tex, OP.tex, OAD.tex, DFD.tex \\
\textbf{Stelle:} Verschiedene Abschnitte \\
\textbf{Problem:}
Die Unterscheidung zwischen **Definition** (terminologisch stabil) und **Modell** (ersetzbar) ist in **MKO.tex** **konzeptionell klar**, aber die **praktische Anwendung** in den Modulen ist nicht immer konsistent.
\begin{itemize}
    \item In **OP.tex** wird das **Fragedifferential** als Definition eingeführt, aber die **Formalisierung als KL-Divergenz** als Modell.
    \item In **OAD.tex** wird **"Fixpunkt"** als Modell verwendet, aber nicht explizit als solches gekennzeichnet.
    \item In **DFD.tex** wird die **KL-Divergenz** als "Heuristik" bezeichnet, aber nicht klar, ob sie als **Modell** oder **Definition** zu verstehen ist.
\end{itemize}
\textbf{Warum bedeutsam:}
Wenn diese Unterscheidung nicht **durchgängig** angewendet wird, könnte dies zu **Kategorienfehlern** führen (z. B. wenn ein Modell fälschlich als Definition behandelt wird). Dies untergräbt die **Kontaminationskontrolle** (MKO).

---

\subsection*{[B-02] Fehlende explizite Verweise auf die Hierarchie der epistemischen Sicherheit in OP.tex}
\textbf{Dokument:} OP.tex \\
\textbf{Stelle:} Verschiedene Axiome und Propositionen \\
\textbf{Problem:}
**MKO** definiert eine **Hierarchie der epistemischen Sicherheit** (Axiom → Proposition → Postulat → Definition/Modell → Bemerkung → Analogie → Hypothese → Offene Forschungsfrage). **OP.tex** verwendet diese Hierarchie **implizit**, aber nicht **explizit als Verweis**.
\begin{itemize}
    \item Axiom 3 (Denk-Monismus) ist ein **Postulat (P)**, aber es wird nicht klar, wie es sich in die **Hierarchie** einordnet.
    \item Die **Minimaldefinition des Denkens** (Axiom 3) wird nicht als **Bemerkung (MKO Ebene 2)** oder **Analogie (MKO Ebene 4)** gekennzeichnet.
\end{itemize}
\textbf{Warum bedeutsam:}
Die Hierarchie ist ein **Werkzeug zur Bewertung der Robustheit** von Aussagen. Fehlende Verweise schwächen die **Nachvollziehbarkeit** des Systems und könnten zu **Überschätzung der epistemischen Gewissheit** führen.

---
\subsection*{[B-03] Unklare Modulabhängigkeiten in DPR.tex}
\textbf{Dokument:} DPR.tex \\
\textbf{Stelle:} Abschnitt "Konzeptuelle Deutungen physikalischer Phänomene" \\
\textbf{Problem:}
Die **Projektive Reduktion (DPR)** setzt die **OAD** und **OPP** voraus, aber es gibt **keine explizite Prüfung**, ob die **OAD tatsächlich das liefert, was die DPR voraussetzt**.
\begin{itemize}
    \item **OAD** beschreibt **Attraktoren** als stabile Konfigurationen im OPP.
    \item **DPR** setzt voraus, dass diese Attraktoren **durch Projektionsoperatoren in raumzeitliche Existenz übersetzt** werden können.
    \item **Frage:** Wird in der OAD ausreichend begründet, dass Attraktoren **tatsächlich die Eigenschaften besitzen**, die für die DPR notwendig sind (z. B. **Selbstkonsistenz, relationale Anschlussfähigkeit**)?
\end{itemize}
\textbf{Warum bedeutsam:}
Gemäß **MKO, modulare Koherenz** müssen die **Schnittstellen zwischen den Modulen konsistent** sein. Fehlende Prüfung könnte zu **Inkonsistenzen** führen, wenn die OAD nicht das liefert, was die DPR benötigt.

---
\subsection*{[B-04] Fehlende empirische Anschlusspunkte in DHP.tex}
\textbf{Dokument:} DHP.tex \\
\textbf{Stelle:} Abschnitt "Dekonstruktion des Hard Problems" \\
\textbf{Problem:}
Die **Dekonstruktion des Hard Problems** basiert auf einer **logischen Analyse** der Zombie-Hypothese und der **Symmetrisierung der Beweislast**. Allerdings fehlen **konkrete empirische Anschlusspunkte**, die zeigen, wie der OP **testbar** ist.
\begin{itemize}
    \item Die **Hypothese H2** (DHP.tex) besagt, dass in allen physikalisch-biologischen Systemen mit struktureller Organisation $A$ auch phänomenales Erleben $B$ vorliegt.
    \item **Frage:** Wie könnte diese Hypothese **empirisch überprüft** werden? Gibt es **experimentelle Designs**, die den OP von anderen Theorien (z. B. Materialismus, Panpsychismus) unterscheiden?
\end{itemize}
\textbf{Warum bedeutsam:}
Gemäß **MKO, Falsifizierbarkeit** sollten Module **explizit benennen**, wo sie **empirisch ansetzbar** sind. Fehlende Anschlusspunkte schwächen die **wissenschaftliche Glaubwürdigkeit** des OP.

---
\subsection*{[B-05] Unklare Rolle des Fragedifferentials in EG-1.tex}
\textbf{Dokument:} EG-1.tex \\
\textbf{Stelle:} Abschnitt "Das Fragedifferential als Individuationsprinzip" \\
\textbf{Problem:}
In **EG-1.tex** wird das **Fragedifferential** als **Individuationsprinzip** eingeführt, aber seine **Rolle im Verhältnis zum Explanatory Gap** bleibt unklar.
\begin{itemize}
    \item **EG-1.tex** behauptet, das Fragedifferential **relokiere** den Explanatory Gap.
    \item **DFD.tex** definiert das Fragedifferential als **ontologisches Konstitutionsprinzip**.
    \item **Frage:** Wie genau **löst** oder **relokiert** das Fragedifferential den Explanatory Gap? Wird der Gap **reduziert**, **verschoben** oder **neu interpretiert**?
\end{itemize}
\textbf{Warum bedeutsam:}
Der **Explanatory Gap** ist das **zentrale Problem** des OP. Unklare Rollenzuweisung könnte zu **Missverständnissen** führen, ob der OP den Gap **tatsächlich adressiert** oder nur **umdefiniert**.

---
%--------------------------------------------------------------
\section*{Minore Befunde [M]}
%--------------------------------------------------------------

\subsection*{[M-01] Terminologische Inkonsistenz: "Denkfeld" vs. "Ontophänomenaler Phasenraum (OPP)"}
\textbf{Dokument:} OP.tex, OPP.tex \\
\textbf{Stelle:} Verschiedene Abschnitte \\
\textbf{Problem:}
In **OP.tex** und **OPP.tex** werden die Begriffe **"Denkfeld"**, **"primordiales Feld"** und **"Ontophänomenaler Phasenraum (OPP)"** **synonym** verwendet, ohne klare Abgrenzung oder Definition der Unterschiede.
\textbf{Warum minor:}
Gemäß **MKO, terminologische Konsistenz** sollten Begriffe **einheitlich** verwendet werden. Inkonsistenzen könnten zu **Verwirrung** führen, insbesondere für neue Leser.

---

\subsection*{[M-02] Fehlende Versionierung in TRANSPARENCY.tex und OP.tex}
\textbf{Dokument:} TRANSPARENCY.tex, OP.tex \\
\textbf{Stelle:} Dokumentenheader \\
\textbf{Problem:}
**MKO.tex** hat eine **klare Versionierung (v1.3)**, während **TRANSPARENCY.tex** und **OP.tex** **keine Versionsnummer** im Header aufweisen.
\textbf{Warum minor:}
Gemäß **MKO, Versionierung** sollten Dokumente **versioniert** sein, um Änderungen nachvollziehbar zu machen.

---
\subsection*{[M-03] Unpräzise Formulierung in DPD.tex: "könnte sich konstituieren"}
\textbf{Dokument:} DPD.tex \\
\textbf{Stelle:} Abschnitt "Das relationale Prinzip der Phänomenalität" \\
\textbf{Problem:}
In der **Definition des phänomenalen Ereignisses** wird der Satz:
\begin{quote}
"Ein phänomenales Ereignis ($\mathcal{P}$) --- das Aufblitzen von Gegebenheit oder Erleben --- könnte sich konzeptuell durch die wechselseitige relationale Verschränkung ($\mathcal{R}$) von mindestens zwei strukturell unterschiedlichen, projizierten Attraktor-Manifestationen ($\mathcal{A}_1$ und $\mathcal{A}_2$) innerhalb des Beobachtungsraums $\mathcal{M}_{\mathcal{O}}$ konstituieren"
\end{quote}
verwendet. Der **Konjunktiv ("könnte")** ist hier **nicht konsistent** mit der **MKO-Sprachebene 1 (Indikativ)**, da es sich um eine **Definition** handelt.
\textbf{Warum minor:}
Gemäß **MKO, Sprachebene 3 (Definitorisch)** sollten Definitionen **präzise und ohne Modalisierungen** formuliert werden.

---
%--------------------------------------------------------------
\section*{Positive Bewertungen}
%--------------------------------------------------------------

\subsection*{[P-01] Klare und transparente Rollenverteilung in TRANSPARENCY.tex}
\textbf{Dokument:} TRANSPARENCY.tex \\
\textbf{Stelle:} Abschnitt "Zusammenarbeit folgt einer klar definierten methodischen Rollenverteilung" \\
\textbf{Stärke:}
Die **drei Rollen** (Konzeptionelle Leitung, Adversarial Reviewer, Formale Ausarbeitung) sind **präzise definiert** und klar voneinander abgegrenzt.
\textbf{Warum solide:}
Verhindert **Missverständnisse** über die Urheberschaft und ermöglicht eine **kritische Reflexion** über den Einsatz von KI. Dies ist ein **Vorbild für Transparenz** in KI-gestützter Forschung.

---

\subsection*{[P-02] Systematische und operationalisierbare R/P-Unterscheidung in MKO.tex}
\textbf{Dokument:} MKO.tex \\
\textbf{Stelle:} Abschnitt "Die R/P-Unterscheidung: Zirkelkontrolle" \\
\textbf{Stärke:}
Die **Rekonstruktiv (R) vs. Postulativ (P)**-Unterscheidung ist **methodisch klar** und **praktisch anwendbar**.
\textbf{Warum solide:}
Ermöglicht eine **kontrollierte Eingrenzung des transzendentalen Zirkels**. Leser können nachvollziehen, welche Aussagen **begründet** und welche **gesetzt** sind. Dies ist ein **zentrales Instrument** für die **Robustheit** des OP.

---
\subsection*{[P-03] Stringente Sprachregel mit fünf Ebenen in MKO.tex}
\textbf{Dokument:} MKO.tex \\
\textbf{Stelle:} Abschnitt "Die Sprachregel: Fünf Ebenen des Aussagenstatus" \\
\textbf{Stärke:}
Die **fünf Sprachebenen** (Indikativ, Konjunktiv, Definitorisch, Modal, Hypothetisch) sind **präzise definiert** und mit **Beispielen** untermauert.
\textbf{Warum solide:}
Verhindert **stille Übergänge** von logischen Schlüssen zu Tatsachenbehauptungen. Leser wissen immer, auf welcher **epistemischen Ebene** sie sich befinden. Dies ist ein **Meilenstein für methodologische Klarheit**.

---
\subsection*{[P-04] Relokation des Explanatory Gaps in EG-1.tex}
\textbf{Dokument:} EG-1.tex \\
\textbf{Stelle:} Abschnitt "Die Position des Ontophänomenalismus" \\
\textbf{Stärke:}
Der OP **relokiert** den Explanatory Gap, anstatt ihn zu **leugnen** oder **umzubenennen**.
\begin{itemize}
    \item Der Gap wird nicht zwischen **Materie und Geist** verortet, sondern zwischen **struktureller Selbstbezüglichkeit** und **phänomenaler Aktualität**.
    \item Die **Zimmer- und Traumanalogie** illustrieren, warum die klassische Fragestellung eine **kategoriale Fehlfrage** sein könnte.
    \item Das **Fragedifferential** wird als **Individuationsprinzip** eingeführt, das den Gap als **konstitutive Grenze** jeder endlichen Perspektive erklärt.
\end{itemize}
\textbf{Warum solide:}
Dies ist eine **ehrliche und innovative** Herangehensweise, die den OP von anderen Theorien abhebt. Statt den Gap zu **ignorieren**, wird er **präzise verortet** und als **notwendige Grenze** akzeptiert.

---
\subsection*{[P-05] Kohärente modulare Architektur}
\textbf{Dokument:} Alle Module (OP, OPP, OAD, DPR, DPD, DFD, EG-1, DHP) \\
\textbf{Stärke:}
Die **modulare Architektur** des OP ist **kohärent** und **gut durchdacht**.
\begin{itemize}
    \item Jedes Modul hat eine **klare Rolle** (z. B. OPP als Fundament, DPR als Brücke zur Raumzeit, DFD als Individuationsprinzip).
    \item Die **Schnittstellen zwischen den Modulen** sind **explizit definiert** (z. B. Übergabe von OAD an DPR).
    \item Die **Abhängigkeiten** sind **logisch strukturiert** (z. B. DPR setzt OPP und OAD voraus).
\end{itemize}
\textbf{Warum solide:}
Dies ermöglicht eine **schrittweise Entwicklung** des OP und erleichtert die **Überprüfung** der einzelnen Komponenten. Die modulare Struktur ist ein **Vorbild für philosophische Systeme**.

---
%--------------------------------------------------------------
\section*{Offene Fragen an das Team}
%--------------------------------------------------------------

\begin{enumerate}
    \item \textbf{Zu TRANSPARENCY.tex:}
    \begin{itemize}
        \item Wie wird sichergestellt, dass **KI-generierte inhaltliche Vorschläge** (z. B. logische Konsolidierungen) **nicht als Ko-Autorschaft** interpretiert werden?
        \item Gibt es ein **Protokoll**, das festhält, welche Teile des OP **ausschließlich vom Autor** stammen und welche **in Zusammenarbeit mit KI** entstanden sind?
    \end{itemize}

    \item \textbf{Zu MKO.tex:}
    \begin{itemize}
        \item Wie wird die **Konsistenz der Definition/Modell-Unterscheidung** über alle Module hinweg **überprüft**?
        \item Gibt es ein **zentrales Glossar**, das alle **Definitionen und Modelle** des OP zusammenfasst?
    \end{itemize}

    \item \textbf{Zu OP.tex:}
    \begin{itemize}
        \item Warum wird der Begriff **"Denkfeld"** nicht durchgängig durch **"Ontophänomenaler Phasenraum (OPP)"** ersetzt (oder umgekehrt)?
        \item Wie wird sichergestellt, dass **alle Interpretationen** (z. B. in Axiom 3) **konsequent im Konjunktiv** formuliert werden?
    \end{itemize}

    \item \textbf{Zu OAD.tex:}
    \begin{itemize}
        \item Warum wird **"Fixpunkt"** nicht explizit als **Analogie (MKO Ebene 4)** gekennzeichnet, sondern als **Definition (MKO Ebene 3)**?
    \end{itemize}

    \item \textbf{Zu DPR.tex:}
    \begin{itemize}
        \item Wie wird sichergestellt, dass die **OAD tatsächlich das liefert, was die DPR voraussetzt** (z. B. Selbstkonsistenz, relationale Anschlussfähigkeit der Attraktoren)?
    \end{itemize}

    \item \textbf{Zu DPD.tex:}
    \begin{itemize}
        \item Wie genau **emergiert** die Subjekt-Objekt-Spaltung aus der **relationalen Asymmetrie**? Gibt es eine **formale Herleitung** oder bleibt dies eine **konzeptuelle Deutung**?
    \end{itemize}

    \item \textbf{Zu DFD.tex:}
    \begin{itemize}
        \item Wie genau **relokiert** das Fragedifferential den Explanatory Gap? Wird der Gap **reduziert**, **verschoben** oder **neu interpretiert**?
    \end{itemize}

    \item \textbf{Zu DHP.tex:}
    \begin{itemize}
        \item Gibt es **konkrete empirische Anschlusspunkte**, die den OP von anderen Theorien (z. B. Materialismus, Panpsychismus) unterscheiden?
    \end{itemize}

    \item \textbf{Zu EG-1.tex:}
    \begin{itemize}
        \item Wie genau **unterscheidet** sich die **kategoriale Fehlfrage** von einer **unbeantwortbaren Frage**? Gibt es Kriterien, um zwischen beiden zu unterscheiden?
    \end{itemize}
\end{enumerate}

%--------------------------------------------------------------
\section*{Priorisierte Empfehlungen}
%--------------------------------------------------------------

Die folgenden drei Befunde sollten **zuerst adressiert** werden, da sie die **methodologische Stringenz** und **Glaubwürdigkeit** des OP direkt betreffen:

\begin{enumerate}
    \item \textbf{[K-01] (TRANSPARENCY.tex):} Klärung der **KI-Rollen** (Adversarial Review vs. Ko-Autorschaft) und **explizite Kennzeichnung** der Urheberschaft.
    \item \textbf{[K-02] (OP.tex):} **Explizite Kennzeichnung der R/P-Unterscheidung** für alle Axiome und Propositionen.
    \item \textbf{[K-03] (OP.tex):} **Korrektur der Sprachebenen** (Indikativ vs. Konjunktiv) für alle Interpretationen.
\end{enumerate}

Anschließend sollten die **bedeutsamen Befunde [B]** angegangen werden, insbesondere:
\begin{enumerate}
    \setcounter{enumi}{3}
    \item \textbf{[B-01] (MKO.tex, OP.tex, OAD.tex, DFD.tex):} Klärung der **Definition/Modell-Unterscheidung** und deren **konsistente Anwendung** in allen Modulen.
    \item \textbf{[B-03] (DPR.tex):} Prüfung der **Modulabhängigkeiten** (OAD → DPR).
\end{enumerate}

\end{document}
